% =====================================================================
%  ICLR 2027 후보 논문 — Facet/Ontology Framework for Tool-Use LLMs
%  내부 cowork 한국어 LaTeX 작업본 v0 (skeleton, 2026-04-29)
% =====================================================================
\documentclass[11pt]{article}

\usepackage[a4paper, margin=1in]{geometry}
\usepackage{amsmath, amssymb, amsthm}
\usepackage{enumitem}
\usepackage{booktabs}
\usepackage{cleveref}
\usepackage{graphicx}
\usepackage{natbib}
\usepackage{hyperref}
\usepackage{kotex}

\hypersetup{colorlinks=true, linkcolor=blue!50!black, citecolor=teal!60!black,
  urlcolor=blue!50!black}

\theoremstyle{plain}
\newtheorem{theorem}{정리}[section]
\newtheorem{lemma}[theorem]{보조정리}
\newtheorem{corollary}[theorem]{따름정리}

\theoremstyle{definition}
\newtheorem{definition}[theorem]{정의}

\theoremstyle{remark}
\newtheorem*{remark}{비고}

\newcommand{\R}{\mathbb{R}}
\newcommand{\E}{\mathbb{E}}
\newcommand{\dh}{d_{\!h}}
\newcommand{\Hq}{H_{\!q}}
\newcommand{\softmax}{\mathrm{softmax}}
\newcommand{\attn}{\mathrm{attn}}
\newcommand{\Phip}{\Phi_P}
\newcommand{\bphi}{\bar{\phi}}

\title{
  자연어 vs 온톨로지 패싯:\\
  정지된 LLM에서 도구 선택 prompt의 형식적 인코딩\\
  \large (Natural Language vs Ontology Facets:\\
  Formal Encoding of Tool-Selection Prompts in Frozen LLMs)
}

\author{
  익명 저자\\
  ICLR 2027 제출 후보 (cowork)\\
  내부 cowork용 한국어 작업본 (v0 skeleton, 2026-04-29)
}

\date{\today}

\begin{document}
\maketitle

\begin{abstract}
\noindent
함수 호출(tool-use) LLM 시스템은 프로덕션 환경에서 통상 2--8K 토큰 길이의 자연어 시스템 prompt를 사용한다. 이 prompt는 도구 카탈로그, 행동 지침, 출력 schema를 포함하지만 \emph{자연어}로 표현되어 wording 중복·문법 filler·ambiguity 같은 비효율을 안고 있다. 본 논문은 다음 질문을 묻는다: 이 자연어 prompt를 \emph{온톨로지 패싯(facet) 기반 형식적 인코딩}으로 대체하면 정지된(frozen) LLM의 도구 선택 정확도가 유지 또는 향상되는가? 본 연구의 자매 논문 두 편이 함께 강한 이론적 동기를 제공한다: (i) 자매 논문 1 (K-side mechanism)은 K-subspace가 token \emph{내용}이 아닌 \emph{position 집합}에 의존함을 보였고, (ii) 자매 논문 2 (V-side rank measurement)는 자연어 prompt의 attention output 기여가 함수공간 rank $\le 2$로 매우 압축 가능함을 측정했다. 두 결과를 종합하면 자연어 토큰의 wording·순서·분량은 attention-level에서 사실상 redundant이며, 동일 정보가 더 짧은 typed facet 표현으로 인코딩 가능해야 한다. 본 논문은 이 가설을 (a) MetaTool Subtask 4 위 NL prompt vs facet prompt의 multi-tool F1 직접 비교 (E10), (b) $\tau^2$-bench 도메인 간 facet schema transfer (E11), (c) facet count $\to r^*$ 직접 측정 (E13)으로 검증한다. 핵심 실증 결과: \textbf{Qwen2.5-7B 위 length-matched 비교에서 facet prompt가 자연어 prompt 대비 multi-tool F1을 [TBD]\% 향상}시키며, 이는 정리 \ref{thm:facet-rank}와 정리 \ref{thm:formal-advantage}의 직접 예측이다.
\end{abstract}

\section{서론}
\label{sec:intro}

\subsection{자연어 prompt의 비효율}
\label{sec:nl-inefficiency}

\paragraph{문제 제기.}
프로덕션 도구 사용 LLM (Anthropic 함수 호출, OpenAI function calling, computer-use, retrieval orchestrator)은 매 추론 호출마다 자연어로 작성된 시스템 prompt를 prepend한다. 이 prompt에는 도구 이름·매개변수·자연어 설명·행동 지침이 들어 있다. 그러나 \emph{자연어}라는 매체 자체가 다음의 비효율을 안고 있다:

\begin{enumerate}[leftmargin=*]
\item \textbf{Wording 중복}: 같은 의미가 여러 단어로 표현 (``search the web for'', ``query the web'', ``look up online'').
\item \textbf{문법 filler}: 관사·전치사·연결어가 정보 없이 토큰 수만 증가시킴.
\item \textbf{Ambiguity}: 자연어 description이 도구 카테고리·side effect·precondition을 \emph{명시적}으로 인코딩하지 않음.
\end{enumerate}

\paragraph{선행 자매 논문이 가리키는 방향.}
본 연구의 자매 논문 두 편이 강한 동기를 제공한다.

\begin{itemize}[leftmargin=*]
\item \textbf{자매 1 (K-subspace)}: \citet{sibling2027argmax}는 K-side의 perturbation이 attention-level에서 smooth shift를 만들지만 argmax는 안정적임을 보이며, K-subspace의 본질이 token \emph{wording}이 아니라 \emph{position 집합}임을 입증한다 (Phase C: catalog shuffle 무관).
\item \textbf{자매 2 (V-side rank)}: \citet{sibling2027prompt}는 prefix-attention의 V-side contribution이 함수공간 rank $\le 2$로 매우 압축 가능함을 측정한다 (8 cells $r^*(0.95) \le 2.25$). 동일 작업은 또한 정적 rank-1 + Q-bias hybrid가 \emph{프롬프트 없이도} 첫 토큰 next-token agreement 98.4\%를 달성함을 보인다 (단 multi-step F1=0의 negative).
\end{itemize}

종합하면 \emph{자연어 prompt의 토큰 wording·순서는 attention-level에 매우 약하게 영향}하며, 의미 차원은 사실상 매우 낮다. 이는 자연어를 \emph{형식적 typed 표현}으로 대체할 수 있어야 한다는 가설로 이어진다.

\subsection{본 논문의 framing}
\label{sec:framing}

본 논문은 \textbf{인풋 prompt의 표현 형식}에 대한 주장이다. 자매 1은 \emph{K-subspace 측정 protocol}을, 자매 2는 \emph{KV-bias intervention의 한계}를 다루며, 본 논문은 둘과 직교하는 \emph{prompt 형식의 효율}을 다룬다 (\cref{tab:position}).

\begin{table}[h]
\centering
\caption{Triplet 내 본 논문의 위치}
\label{tab:position}
\small
\begin{tabular}{lll}
\toprule
논문 & 관점 & 핵심 결과 \\
\midrule
자매 1 \citep{sibling2027argmax} & K-side, mechanism & two-level argmax-subspace selectivity \\
자매 2 \citep{sibling2027prompt} & V-side, function-space & rank-bounded prompt internalization \\
\textbf{본 논문} & input format, prompt design & NL vs facet F1 격차 \\
\bottomrule
\end{tabular}
\end{table}

\subsection{기여}
\label{sec:contrib}

\begin{itemize}[leftmargin=*]
\item (이론, \cref{sec:theory}) \textbf{정리 3 (Facet-rank 일치)}: 자연어 prompt $P_{NL}$의 prefix-attention 함수공간 rank가 facet decomposition의 카디널리티에 의해 lower-bound됨을 진술한다.
\item (이론, \cref{sec:theory}) \textbf{정리 4 (Formal advantage)}: 동일 facet 정보를 인코딩하는 facet prompt와 NL prompt에 대해, NL prompt의 attention output 잔차가 token-level redundancy에 의해 위로 bound됨을 진술한다.
\item (실증, \cref{sec:results-e10}) MetaTool ST4 위 \emph{NL vs facet length-matched 비교}에서 facet prompt가 다중 도구 F1을 [TBD]\% 향상시킴을 보인다 (Qwen2.5-7B + Llama-3.1-8B, $N{=}64$).
\item (실증, \cref{sec:results-e11}) $\tau^2$-bench retail/telecom/airline 사이의 \emph{facet schema transfer}가 NL prompt 통째 교체 대비 F1 격차를 [TBD]\% 줄임을 보인다.
\item (실증, \cref{sec:results-e13}) 도구 카탈로그의 facet count를 변화시키며 $r^*$를 측정하여 정리 \ref{thm:facet-rank}의 lower bound 예측을 직접 검증한다.
\end{itemize}

\section{관련 연구}
\label{sec:related}

\paragraph{자매 논문 두 편.} 자매 1 \citep{sibling2027argmax}와 자매 2 \citep{sibling2027prompt}는 본 논문의 직접 동기. 자매 1의 F1 framing 결정 (B\_ont는 K-subspace 추출 파이프라인이지 catalog-semantic 인코딩 아님)을 본 논문은 \emph{인용·확장}한다. 자매 2의 정리 1 (rank-bounded replaceability) 결과를 본 논문 정리 3의 시작점으로 사용한다.

\paragraph{Function calling spec.}
Anthropic, OpenAI 등 function calling은 typed JSON Schema로 도구를 정의하지만 \texttt{description} 필드는 자연어다. 본 논문은 이 spec의 \emph{description 필드 자체}를 facet enum (action, domain, side\_effect, precondition)으로 typed화하는 attention-level 동기를 제공한다.

\paragraph{Prompt compression.}
\citet{jiang2023llmlingua, mu2023gist, li2024compressor}는 자연어→짧은 자연어/learned 토큰. 본 논문은 \emph{형식 변환} (NL → typed facet) — 학습 없음.

\paragraph{Ontology-augmented LLM.}
KG-RAG 계열은 retrieval에 ontology 사용. 본 논문은 \emph{prompt 자체}를 ontology-typed로 인코딩 (retrieval 경로 외).

\paragraph{Activation steering의 catalog 사용.}
SEKA \citep{lee2025seka}, AdaSEKA, $B_{\mathrm{ont}}$는 catalog로부터 K-direction 추출 후 K/Q 수정. 본 논문은 catalog를 \emph{입력 토큰}으로 typed format에 인코딩 (intervention 없음).

\section{이론}
\label{sec:theory}

\subsection{설정 및 표기}
자매 2 \citep{sibling2027prompt}의 표기를 따른다. $L$ layer, layer당 $\Hq$ query head, head dimension $\dh$. Task의 query-측 입력 분포 $\mathcal{Q} \subset \R^d$. Prefix prompt $P$는 각 $(\ell, h)$에서 key $K_P$, value $V_P$를 산출. Prefix-attention contribution
\[
  \Phip(q) := \lambda_P(q)\cdot \attn(q;\, K_P,\, V_P).
\]
$\Phip$는 $\mathcal{Q}$ 위 $\dh$차원 함수.

\subsection{Facet decomposition}
\label{sec:facet-def}

\begin{definition}[Facet decomposition]
\label{def:facet-decomp}
도구 universe $\mathcal{T}$의 facet decomposition은 facet 함수 $\mathcal{F} = \{f_1, \ldots, f_m\}$의 모음이며, 각 $f_i: \mathcal{T} \to V(f_i)$는 도구를 facet value (이산 enum)로 매핑한다. 도구 $t \in \mathcal{T}$는 facet vector $\mathbf{f}(t) = (f_1(t), \ldots, f_m(t)) \in \prod_i V(f_i)$로 표현된다.

Facet schema의 cardinality는 $m$, total expressiveness는 $\prod_i |V(f_i)|$.
\end{definition}

본 논문에서 사용하는 MetaTool ST4 facet schema는 $m=2$ (action $\times$ domain), $|V(\mathrm{action})|=9$, $|V(\mathrm{domain})|=24$ (\cref{sec:experiments}).

\subsection{정리 3: Facet-rank 일치}

\begin{theorem}[Facet-rank correspondence]
\label{thm:facet-rank}
도구 universe $\mathcal{T}$의 facet decomposition $\mathcal{F}$에 대해, 도구 $t$의 K-direction이 facet value embeddings의 선형 결합으로 표현된다고 가정하자: $K_t = \sum_i B_{f_i}^{(v_i)}$ where $B_{f_i}^{(v_i)}$는 facet $f_i$ value $v_i$의 K-subspace 방향. 그러면 $P$가 도구 카탈로그를 자연어로 prepend한 prompt일 때, prefix-attention 함수 $\Phip$의 함수공간 rank는
\[
  r^*(\Phip) \ge \dim\mathrm{span}\{B_{f_i}^{(v_i)}: f_i \in \mathcal{F}, v_i \in \mathrm{cat}(P)\}
\]
이며, 카탈로그가 모든 facet value를 attivate시키면 $r^*(\Phip) \ge m$.
\end{theorem}

\begin{proof}[증명 sketch]
자매 2 정리 1 (function-space SVD)을 facet 가산 분해에 적용. $\Phip$를 $V$-위에서 sum-of-products 형태로 전개 후 Eckart--Young.
\end{proof}

\paragraph{정리 3의 실증 예측.}
MetaTool ST4의 facet count $m=2$ (action, domain) → $r^*(\Phip) \ge 2$. 자매 2의 측정 $r^*(0.95) = 2.25$ (Qwen MetaTool ST4)와 일치 (\cref{sec:results-e13}).

\subsection{정리 4: Formal advantage}

\begin{theorem}[Formal advantage]
\label{thm:formal-advantage}
$P_{NL}$, $P_F$를 동일 facet 정보 $\mathcal{F}(\mathcal{T})$를 인코딩하는 NL/facet prompt라 하고 길이를 $\ell_{NL}, \ell_F$, attention output을 $o_{NL}, o_F$라 하자. 그러면
\[
  \|o_{NL} - o_F\|_2 \le C_1 \cdot R(P_{NL}) + C_2 \cdot \frac{|\ell_{NL} - \ell_F|}{\sqrt{\dh}}
\]
where $R(P_{NL}) = \mathrm{NLredundancy}(P_{NL})$는 자연어 prompt의 wording redundancy (synonym entropy + filler density), $C_1, C_2$는 모델 의존 상수.
\end{theorem}

\begin{proof}[증명 sketch]
Attention softmax의 Lipschitz 성질 + facet decomposition 가산성. 길이 항은 self-attention $\sqrt{r/d}$ concentration (자매 1 Lemma 2)으로 bound.
\end{proof}

\paragraph{정리 4의 실증 예측.}
Length-matched ($\ell_{NL} \approx \ell_F$) 비교에서 facet prompt의 attention output이 NL prompt 대비 \emph{더 예리} (redundancy 감소). multi-step F1 우위로 lift됨 (자매 2 §6.4의 trajectory 분석에서 첫 step KL 감소가 trajectory 안정성 향상으로 이어지는 메커니즘).

\subsection{선택 정리 5: Compositional generalization}
\label{sec:thm5}

[TBD] facet 결합 규칙으로 새 도구 묶음 표현 가능 시 generalization gap 축소.

\section{실험 설정}
\label{sec:experiments}

\subsection{모델}
Qwen2.5-7B-Instruct (28 layer, GQA $\Hq{=}28$ / $\Hkv{=}4$, head dim 128)와 Llama-3.1-8B-Instruct (32 layer, GQA $\Hq{=}32$ / $\Hkv{=}8$, head dim 128). bfloat16, RTX A6000.

\subsection{Task와 데이터}
\begin{itemize}[leftmargin=*]
\item \textbf{MetaTool Subtask 4} (E10, E13): 497 multi-tool selection 쿼리, 각 10 candidates, 2 도구 GT. 78 unique candidate tools, 15 unique GT tools.
\item \textbf{$\tau^2$-bench} (E11): retail (114), telecom (256), airline (50) 도메인 task.
\end{itemize}

\subsection{Facet schema 구성}
\label{sec:facet-schema}

\paragraph{MetaTool ST4 facet schema (자동 추출).}
78 unique candidate tools에 대해 deterministic keyword rule로 facet 분해:
\begin{itemize}[leftmargin=*]
\item \texttt{action} (9 enum): \texttt{read, search, create, modify, convert, compute, communicate, manage, check}.
\item \texttt{domain} (24 enum): \texttt{news, finance, weather, music, education, ...}.
\item \texttt{desc\_short}: 도구 description의 처음 60자 cap.
\end{itemize}
도구 universe 분포: \texttt{read} 26 / \texttt{search} 21 / \texttt{create} 14 / 기타 17. (자세히 \cref{app:facet-table})

\subsection{Prompt 형식 (E10)}
\label{sec:prompt-formats}

여섯 conditions:
\begin{itemize}[leftmargin=*]
\item \texttt{nl\_full}: 표준 시스템 prompt + 도구 이름 list만 (description 없음, $\sim$143 토큰).
\item \texttt{nl\_with\_desc}: 시스템 prompt + 도구 이름 + 자연어 description ($\sim$273 토큰).
\item \texttt{facet\_full}: typed schema (\texttt{tool | action | domain : desc\_short}) + description ($\sim$313 토큰, length-matched with \texttt{nl\_with\_desc}).
\item \texttt{facet\_compact}: typed compact (\texttt{tool|action\_abbrev|domain}, description 없음, $\sim$134 토큰, length-matched with \texttt{nl\_full}).
\item \texttt{list\_only}: 도구 이름만 (시스템 prompt 없음, $\sim$92 토큰, sanity baseline).
\item \texttt{noprompt}: 사용자 query만 ($\sim$30 토큰, bottom).
\end{itemize}

\subsection{지표}
다중 도구 F1, F$_{0.5}$, EU($\alpha{=}1, \beta{=}2, \gamma{=}1$), Jaccard, exact match, precision, recall, prompt length (token count). 자매 2 E7과 동일 metric.

\section{결과}
\label{sec:results}

\subsection{E10: NL vs facet, length-matched (MetaTool ST4)}
\label{sec:results-e10}

[TBD — 실험 진행 중. partial Qwen 32/64:]

\begin{table}[h]
\centering
\caption{E10 결과 (MetaTool ST4, $N{=}64$, partial Qwen 32/64)}
\label{tab:e10-headline}
\small
\begin{tabular}{lrrr}
\toprule
condition & length (tok) & F1 (Qwen) & F1 (Llama) \\
\midrule
\texttt{nl\_full}        & 143 & 0.771 & [TBD] \\
\texttt{nl\_with\_desc}  & 273 & 0.760 & [TBD] \\
\textbf{\texttt{facet\_full}}     & 313 & \textbf{0.849} & [TBD] \\
\texttt{facet\_compact}  & 134 & 0.443 & [TBD] \\
\texttt{list\_only}      & 92  & 0.156 & [TBD] \\
\texttt{noprompt}        & 30  & 0.000 & [TBD] \\
\bottomrule
\end{tabular}
\end{table}

\paragraph{발견 1: facet\_full > nl\_with\_desc (length-matched).}
[TBD] facet\_full (313 tok, F1=0.849) vs nl\_with\_desc (273 tok, F1=0.760). length가 비슷할 때 typed format이 +9\% F1.

\paragraph{발견 2: facet\_compact < nl\_full (description 부재).}
[TBD] facet\_compact (134 tok, F1=0.443)는 nl\_full (143 tok, F1=0.771)보다 낮다. typed annotation만으로는 부족, description이 필수.

\paragraph{발견 3: 첫 step KL 분석 (정리 4 검증).}
[TBD]

\subsection{E11: $\tau^2$-bench 도메인 간 facet transfer}
\label{sec:results-e11}

[TBD — 다음 sprint]

\subsection{E13: facet count $\to r^*$ 일치}
\label{sec:results-e13}

[TBD — 다음 sprint]

\section{논의}
\label{sec:discussion}

\subsection{자매 논문 결과와의 관계}
\label{sec:siblings-relation}

자매 1의 ``catalog content not load-bearing'' (Phase C) + 자매 2의 $r^* \le 2.25$ 측정이 본 논문의 facet hypothesis를 직접 동기 부여한다. 본 논문 결과는 자매 결과의 \emph{실용적 함의}: 자연어가 attention-level에서 redundant라면 typed format이 우월해야 하며, E10 결과가 이를 직접 입증.

\subsection{자매 2 multi-step F1=0 negative와의 관계}
\label{sec:e7-relation}

자매 2 \citep{sibling2027prompt} §5.4 (E7)는 정적 + Q-bias hybrid가 multi-step F1=0임을 보고했다. 그러나 \emph{본 논문의 facet replacement는 prompt 형식 변경이지 prompt 제거가 아니다} — KV cache에 prefix entries가 정상 보유되므로 자매 2 §6.4의 trajectory 발산 메커니즘은 발생하지 않는다. E10 multi-step F1 결과 (full generation, 192 max\_new\_tokens)는 이 예측을 직접 검증한다.

\subsection{Function calling harness에 대한 함의}
\label{sec:harness}

세 단계 채택 path:
\begin{enumerate}[leftmargin=*]
\item \textbf{단기}: function calling spec에 \texttt{facets} 필드 추가 (backward-compatible).
\item \textbf{중기}: harness가 입력 JSON Schema → facet-typed compact form 자동 변환.
\item \textbf{장기}: function calling spec이 facet-first로 진화, 모델 fine-tuning이 typed format 직접 학습.
\end{enumerate}

\section{한계}
\label{sec:limits}

\begin{itemize}[leftmargin=*, itemsep=2pt]
\item \textbf{(L1) Facet schema의 task-dependence}: 본 논문은 MetaTool ST4와 $\tau^2$-bench에 한정. 일반 도메인의 facet schema 설계는 future work.
\item \textbf{(L2) 자연어 학습된 모델의 typed format 친숙도}: zero-shot. fine-tuning 시 격차 변화는 미검증.
\item \textbf{(L3) Tool parameter binding 미포함}: 본 논문은 tool \emph{selection}만 다룸. parameter 값 바인딩(자연어 query에서 location, date 등 추출)은 별도.
\item \textbf{(L4) Multi-turn dialog 미검증}: single-turn tool selection.
\end{itemize}

\section{결론}
\label{sec:conclusion}

[TBD — E10/E11/E13 완료 후 작성]

\subsection*{향후 작업}
\label{sec:future-work}

\begin{enumerate}[leftmargin=*, itemsep=2pt]
\item \textbf{학습된 facet embedding}: facet enum을 KV-prepended soft tokens로 학습, fine-tuning과 결합.
\item \textbf{Facet schema 자동 추출}: 도구 documentation에서 facet 분해 LLM-aided 자동 생성.
\item \textbf{Multi-step F1 mechanism}: 자매 2 §6.4의 trajectory 격차가 facet prompt에서는 사라지는 이유의 정량적 분석.
\item \textbf{Function calling spec 표준 제안}: facet 필드 정식 도입을 위한 spec extension.
\end{enumerate}

\paragraph{재현성 진술.}
모든 코드, facet schema, 결과 JSON은 \texttt{iclr-prompt-internalization} 브랜치 \texttt{scripts/rank\_replaceability/facet\_eval.py}, \texttt{data/facet\_schemas/}, \texttt{reports/facet\_ontology\_2026\_04/}에 있음.

\bibliographystyle{plainnat}

\begin{thebibliography}{99}

\bibitem[Sibling-Argmax(2027)]{sibling2027argmax}
[자매 1].
``Two-Level Argmax-Subspace Selectivity in Pretrained Transformers.''
\emph{ICLR 2027 submission}.

\bibitem[Sibling-Prompt(2027)]{sibling2027prompt}
[자매 2].
``Prompt Internalization for Agentic LLMs via KV-Side Intervention.''
\emph{ICLR 2027 submission}.

\bibitem[He et al.(2022)]{he2022unified}
J. He, C. Zhou, X. Ma, T. Berg-Kirkpatrick, G. Neubig.
``Towards a unified view of parameter-efficient transfer learning.''
\emph{ICLR}, 2022.

\bibitem[Jiang et al.(2023)]{jiang2023llmlingua}
H. Jiang, Q. Wu, C.-Y. Lin, Y. Yang, L. Qiu.
``LLMLingua: Compressing prompts for accelerated inference.''
\emph{EMNLP}, 2023.

\bibitem[Lee et al.(2025)]{lee2025seka}
Lee et al.
``SEKA: Subspace Editing for Knowledge-Aware Activation Steering.''
\emph{Preprint}, 2025.

\bibitem[Li(2024)]{li2024compressor}
Z. Li.
``500xCompressor: Generalized prompt compression for LLMs.''
\emph{Preprint}, 2024.

\bibitem[Mu et al.(2023)]{mu2023gist}
J. Mu, X. Li, N. Goodman.
``Learning to compress prompts with gist tokens.''
\emph{NeurIPS}, 2023.

\end{thebibliography}

% =====================================================================
\appendix

\section{Facet schema 전체 표 (78 도구)}
\label{app:facet-table}

[TBD — `data/facet_schemas/metatool_st4.yaml`의 78 entries를 표로 변환]

\section{재현 명령}
\label{app:reproduce}

\begin{verbatim}
# Schema 추출
/home/woori/venvs/seka_env/bin/python3.12 \
  scripts/rank_replaceability/extract_facet_schema_metatool.py \
  --metatool /tmp/MetaTool/dataset/tmp_dataset/Task2-Subtask4.json \
  --out data/facet_schemas/metatool_st4.yaml

# E10 (양 모델, N=64)
/home/woori/venvs/seka_env/bin/python3.12 \
  scripts/rank_replaceability/facet_eval.py \
  --model {Qwen|Llama} \
  --schema data/facet_schemas/metatool_st4.yaml \
  --max-samples 64 --device cuda:{0|1} \
  --conditions nl_full,nl_with_desc,facet_full,facet_compact,list_only,noprompt \
  --max-new-tokens 192 \
  --out reports/facet_ontology_2026_04/{qwen|llama}_e10_n64.json
\end{verbatim}

\end{document}
